\lecture{34}{Капацитированное экономическое диспетчирование}

\sourcevideo
    {https://www.youtube.com/playlist?list=PLOzRYVm0a65fhKDS8cYIC7YCuVGJ4FOJx}
    {Week 9: Lecture 34: Capacitated EDP}

Учитываются физические ограничения генераторов
\[
    P_i^{\min}\le P_i\le P_i^{\max}.
\]
Они изменяют условия оптимальности: предельные стоимости свободных
генераторов совпадают, а генераторы на границах удовлетворяют
односторонним условиям KKT.

\section{Задача и условия KKT}

Пусть
\[
    C_i(P_i)=\alpha_iP_i^2+\beta_iP_i+\gamma_i,
    \qquad \alpha_i>0.
\]
Рассматривается выпуклая задача
\[
\begin{aligned}
    \min_{P_1,\ldots,P_N}\quad
    &\sum_{i=1}^{N}C_i(P_i),
    \\
    \text{при условиях}\quad
    &\sum_{i=1}^{N}P_i=P_{\mathrm D},
    \\
    &P_i^{\min}\le P_i\le P_i^{\max}.
\end{aligned}
\]
Допустимость требует
\[
    \sum_{i=1}^{N}P_i^{\min}
    \le P_{\mathrm D}\le
    \sum_{i=1}^{N}P_i^{\max}.
\]
Запишем ограничения мощности как
\[
    P_i-P_i^{\max}\le0,
    \qquad
    P_i^{\min}-P_i\le0.
\]
Для баланса введём множитель \(\lambda\in\mathbb R\), а для верхнего
и нижнего ограничений --- множители \(r_i,s_i\ge0\). Функция Лагранжа
равна
\[
\begin{aligned}
    \mathcal L(P,\lambda,r,s)
    ={}&
    \sum_{i=1}^{N}C_i(P_i)
    +\lambda\left(P_{\mathrm D}-\sum_{i=1}^{N}P_i\right)
    \\
    &+
    \sum_{i=1}^{N}r_i(P_i-P_i^{\max})
    +
    \sum_{i=1}^{N}s_i(P_i^{\min}-P_i).
\end{aligned}
\]
Условия KKT имеют вид
\[
    2\alpha_iP_i^\star+\beta_i
    -\lambda^\star+r_i^\star-s_i^\star=0,
\]
\[
    \sum_{i=1}^{N}P_i^\star=P_{\mathrm D},
    \qquad
    P_i^{\min}\le P_i^\star\le P_i^{\max},
\]
\[
    r_i^\star,s_i^\star\ge0,
    \qquad
    r_i^\star(P_i^\star-P_i^{\max})=0,
    \qquad
    s_i^\star(P_i^{\min}-P_i^\star)=0.
\]
Следовательно, для свободного генератора
\[
    P_i^{\min}<P_i^\star<P_i^{\max}
    \quad\Longrightarrow\quad
    \lambda^\star=2\alpha_iP_i^\star+\beta_i,
\]
на верхней границе
\[
    P_i^\star=P_i^{\max}
    \quad\Longrightarrow\quad
    \lambda^\star\ge2\alpha_iP_i^{\max}+\beta_i,
\]
а на нижней границе
\[
    P_i^\star=P_i^{\min}
    \quad\Longrightarrow\quad
    \lambda^\star\le2\alpha_iP_i^{\min}+\beta_i.
\]

\section{Проекция мощности и скалярное уравнение}

Для фиксированного \(\lambda\) локальная оптимальная реакция
генератора равна
\[
    P_i(\lambda)
    =
    \min\left\{
        P_i^{\max},
        \max\left\{
            P_i^{\min},
            \frac{\lambda-\beta_i}{2\alpha_i}
        \right\}
    \right\}.
\]
Тем самым некапацитированное значение проецируется на интервал
\([P_i^{\min},P_i^{\max}]\). Оптимальный множитель определяется
уравнением
\[
    \sum_{i=1}^{N}P_i(\lambda^\star)=P_{\mathrm D}.
\]
Левая часть непрерывна и не убывает по \(\lambda\). При наличии хотя
бы одного свободного генератора она строго возрастает в окрестности
решения, поэтому \(\lambda^\star\) определяется однозначно. Если все
генераторы находятся на границах, мощности уже фиксированы, а сам
множитель может быть неоднозначен на интервале, совместимом с KKT.

Уравнение можно решать централизованно бисекцией. В распределённой
схеме вместо прямого поиска корня используется активное множество и
несколько раундов среднего консенсуса.

\section{Активное множество и формула множителя}

Пусть \(\Theta\) --- множество генераторов, работающих на границах, и
\[
    \widehat P_i\in\{P_i^{\min},P_i^{\max}\},
    \qquad i\in\Theta.
\]
Для свободных генераторов
\[
    P_i^\star
    =
    \frac{\lambda^\star-\beta_i}{2\alpha_i},
    \qquad i\notin\Theta.
\]
Из баланса мощности получаем
\[
    \lambda^\star
    =
    \frac{
        P_{\mathrm D}
        -\displaystyle\sum_{i\in\Theta}\widehat P_i
        +\displaystyle\sum_{i\notin\Theta}
        \frac{\beta_i}{2\alpha_i}
    }{
        \displaystyle\sum_{i\notin\Theta}
        \frac1{2\alpha_i}
    }.
\]

Обозначим через \(\lambda_{\mathrm u}^\star\) множитель
некапацитированной задачи:
\[
    \lambda_{\mathrm u}^\star
    =
    \frac{
        P_{\mathrm D}
        +\displaystyle\sum_{i=1}^{N}
        \frac{\beta_i}{2\alpha_i}
    }{
        \displaystyle\sum_{i=1}^{N}
        \frac1{2\alpha_i}
    },
\]
\[
    P_{i,\mathrm u}^\star
    =
    \frac{\lambda_{\mathrm u}^\star-\beta_i}{2\alpha_i}.
\]
Тогда для фиксированного активного множества
\[
    \lambda^\star
    =
    \lambda_{\mathrm u}^\star
    +
    \frac{
        \displaystyle\sum_{i\in\Theta}
        \left(P_{i,\mathrm u}^\star-\widehat P_i\right)
    }{
        \displaystyle\sum_{i\notin\Theta}
        \frac1{2\alpha_i}
    }.
\]
Числитель показывает суммарное изменение мощности из-за насыщения, а
знаменатель --- чувствительность свободных генераторов к изменению
общего множителя.

\section{Распределённое вычисление поправки}

Для заданного множества \(\Theta\) агент \(i\) инициализирует
\[
    y_i(0)
    =
    \begin{cases}
        P_{i,\mathrm u}^\star-\widehat P_i,
        &i\in\Theta,\\[1mm]
        0,
        &i\notin\Theta,
    \end{cases}
\]
и
\[
    z_i(0)
    =
    \begin{cases}
        0,
        &i\in\Theta,\\[1mm]
        \dfrac1{2\alpha_i},
        &i\notin\Theta.
    \end{cases}
\]
Два параллельных алгоритма среднего консенсуса дают всем агентам
\[
    \overline y
    =
    \frac1N
    \sum_{i\in\Theta}
    \left(P_{i,\mathrm u}^\star-\widehat P_i\right),
    \qquad
    \overline z
    =
    \frac1N
    \sum_{i\notin\Theta}
    \frac1{2\alpha_i}.
\]
Множители \(1/N\) сокращаются, поэтому каждый агент вычисляет
\[
    \lambda^\star
    =
    \lambda_{\mathrm u}^\star
    +
    \frac{\overline y}{\overline z}.
\]
После этого свободные мощности пересчитываются по формуле
\[
    P_i
    =
    \frac{\lambda^\star-\beta_i}{2\alpha_i},
    \qquad i\notin\Theta,
\]
а активные остаются на назначенных границах. Коэффициенты
\(\alpha_i,\beta_i,\gamma_i\) напрямую не передаются, однако обмен
состояниями \(y_i,z_i\) не является криптографической гарантией
приватности.

\section{Активно-множественная процедура}

Сначала распределённо решается некапацитированная задача и строятся
\(\lambda_{\mathrm u}^\star\) и \(P_{i,\mathrm u}^\star\). Затем при
текущем \(\lambda\) каждый свободный генератор вычисляет
\[
    P_i=\frac{\lambda-\beta_i}{2\alpha_i}.
\]
Нарушители образуют множество
\[
\begin{aligned}
    \Omega
    =
    \bigl\{
        i\notin\Theta:
        P_i<P_i^{\min}
        \text{ или }
        P_i>P_i^{\max}
    \bigr\}.
\end{aligned}
\]
Для \(i\in\Omega\) мощность фиксируется на ближайшей нарушенной
границе:
\[
    \widehat P_i
    =
    \begin{cases}
        P_i^{\min},&P_i<P_i^{\min},\\[1mm]
        P_i^{\max},&P_i>P_i^{\max}.
    \end{cases}
\]
После обновления \(\Theta\) через два консенсуса пересчитывается
\(\lambda\), затем пересчитываются мощности свободных генераторов.
Процедура завершается, когда нарушителей нет и одновременно выполнены
односторонние условия KKT для активных индексов:
\[
    C_i'(P_i^{\max})\le\lambda^\star
    \quad\text{на верхней границе},
\]
\[
    C_i'(P_i^{\min})\ge\lambda^\star
    \quad\text{на нижней границе}.
\]
Одной проверки допустимости мощностей недостаточно.

В лекции активное множество только увеличивается. Это корректно, если
однажды насыщенный генератор после каждого пересчёта остаётся
активным. В общей задаче при одновременных нарушениях верхних и нижних
границ такое свойство требует отдельного доказательства. Надёжная
реализация должна разрешать удаление индексов из \(\Theta\) либо
проверять конечную точку через проекционное уравнение
\[
    \sum_{i=1}^{N}P_i(\lambda)=P_{\mathrm D}.
\]

\section{Завершение и границы модели}

Если на каждом раунде к корректному активному множеству добавляется
хотя бы один новый индекс и удаление не требуется, число обновлений не
превосходит \(N\). Если некапацитированная задача решается за время
\(T_{\mathrm u}\), а один раунд двух параллельных консенсусов занимает
не более \(T_{\mathrm c}\), то
\[
    T_{\mathrm{total}}
    \le
    T_{\mathrm u}+NT_{\mathrm c}.
\]
Эта оценка не зависит от начальных значений непрерывных консенсусных
переменных, но зависит от числа генераторов и коммуникационного графа.

Рассмотренная модель не учитывает потери мощности, ограничения потоков
по линиям, динамику генераторов, невыпуклые затраты, стоимость запуска
и изменение нагрузки во времени. Точная фиксированно-временная
гарантия относится к непрерывным консенсусным динамикам; после
дискретизации требуется отдельный анализ устойчивости и времени
сходимости.
